% Unit 089 terjemahan Bahasa Indonesia dari chapter7.tex baris 419--610.
% Sumber: Methods of Algebra, Volume 2, commit
% 9a5803ff77dd3257484cb177f851a73770a59dd3 (CC BY 4.0).
% stable-unit-id: o014.aljabr2.chapter7.differential-graded-structures
% Cakupan: Bagian 7.2 lengkap, ``Contoh: Struktur Diferensial Bergradasi'';
% berhenti tepat sebelum Bagian 7.3 pada baris 611.

% segment-id: o014.li2.chapter7.differential-graded-structures.h001
\section{Contoh: Struktur Diferensial Bergradasi}\label{sec:dga-monoidal}
\hypertarget{o014.aljabr2.chapter7.differential-graded-structures}{ }
% segment-id: o014.li2.chapter7.differential-graded-structures.p001
Melanjutkan pembahasan dalam Contoh~\ref{eg:graded-module}, tinjau kategori
abelian monoidal simetris $(\mathcal{A}, \otimes)$ yang memiliki koproduk
terhitung. Dalam bagian ini, kita juga mengandaikan bahwa $\otimes$ eksak
kanan pada setiap peubah. Contoh khasnya tentu saja
$\mathcal{A}=\Bbbk\dcate{Mod}$.

% segment-id: o014.li2.chapter7.differential-graded-structures.p002
Sebelumnya telah dibahas aljabar bergradasi-$\Z$ di atas $\mathcal{A}$, atau,
dengan kata lain, aljabar di atas kategori monoidal
$(\mathcal{A}^{\Z},\otimes)$. Lebih tepatnya, dalam kerangka
Contoh~\ref{eg:graded-module}, hal ini berarti mengambil $I=\Z$ dan struktur
kepang Koszul~\eqref{eqn:Koszul-braiding} yang bersesuaian dengan
$\epsilon(a)=a\bmod\;2$. Selanjutnya, ``bergradasi'' selalu berarti
``bergradasi-$\Z$''. Fokus bagian ini ialah menambahkan struktur
``diferensial''; sebagian pembahasannya bertumpang tindih dengan
\S\ref{sec:multiplicative-SS}.

% segment-id: o014.li2.chapter7.differential-graded-structures.p003
Funktor pergeseran $T:\mathcal{A}^{\Z}\to\mathcal{A}^{\Z}$ didefinisikan
oleh $(TM)^n=M^{n+1}$ dan $(Tf)^n=f^{n+1}$. Dengan demikian,
$\left(\mathcal{A}^{\Z},T\right)$ menjadi kategori abelian dengan pergeseran,
dan objek diferensial di atasnya (Definisi~\ref{def:differential-object})
disingkat menjadi \emph{objek diferensial bergradasi} di atas $\mathcal{A}$.
Sesungguhnya, kategori abelian
$\left(\mathcal{A}^{\Z},T\right)_d$ yang terdiri atas objek-objek diferensial
bergradasi secara kanonik isomorfik dengan kategori kompleks
$\cate{C}(\mathcal{A})$. Fakta ini telah diketahui sejak
\S\ref{sec:additive-cplx}.
\index{objek dg@objek diferensial bergradasi}

% segment-id: o014.li2.chapter7.differential-graded-structures.p004
Karena itu, selanjutnya kita akan menyamakan objek diferensial bergradasi di
atas $\mathcal{A}$ dengan kompleks, tanpa membedakan keduanya.

% segment-id: o014.li2.chapter7.differential-graded-structures.p005
Sekarang berikan struktur kategori monoidal simetris pada
$\cate{C}(\mathcal{A})$. Untuk sembarang kompleks $M$ dan $N$, bifunktor
$\otimes:\mathcal{A}\times\mathcal{A}\to\mathcal{A}$ memberikan bikompleks
$M^\bullet\otimes N^\bullet$; dengan mengambil kompleks total melalui jumlah
langsung, diperoleh
% segment-id: o014.li2.chapter7.differential-graded-structures.q001
\[ M \otimes N := \tot_{\oplus}\left( M^\bullet \otimes N^\bullet \right), \quad (M \otimes N)^n = \bigoplus_{a+b=n} M^a \otimes N^b. \]

% segment-id: o014.li2.chapter7.differential-graded-structures.p006
Dalam contoh $\mathcal{A}=\Bbbk\dcate{Mod}$, diferensial
$d=d_{M\otimes N}$ pada kompleks (atau, secara ekuivalen, pada objek
diferensial) dapat ditulis pada tingkat unsur sebagai
% segment-id: o014.li2.chapter7.differential-graded-structures.q002
\begin{equation*}
	d(x \otimes y) = (d_M x) \otimes y + (-1)^a x \otimes (d_N y), \quad x \in M^a, \; y \in N^b.
\end{equation*}
Definisi-definisi ini membuat struktur anyaman Koszul
\eqref{eqn:Koszul-braiding} memberikan morfisme antarkompleks. Pembaca dapat
memeriksanya langsung atau menurunkannya dari
Proposisi~\ref{prop:double-cplx-swap}.

% segment-id: o014.li2.chapter7.differential-graded-structures.p007
Selain itu, objek satuan $\munit$ dari $(\mathcal{A},\otimes)$, dipandang
sebagai kompleks yang terkonsentrasi pada derajat nol, jelas memberikan objek
satuan bagi kategori monoidal $(\cate{C}(\mathcal{A}),\otimes)$. Fakta berikut
langsung terlihat: untuk setiap
$X=(X^n,d_X^n)_n\in\Obj(\cate{C}(\mathcal{A}))$, berlaku
% segment-id: o014.li2.chapter7.differential-graded-structures.q003
\begin{equation}\label{eqn:unit-map-into-cplx}
	\Hom_{\cate{C}(\mathcal{A})}(\munit, X) \simeq \Hom_{\mathcal{A}}\left( \munit, \Ker\left(d_X^0 \right) \right).
\end{equation}

% segment-id: o014.li2.chapter7.differential-graded-structures.d001
\begin{definition}\label{def:dg-algebra}
	\index{aljabar dg@aljabar diferensial bergradasi}
	\index{aljabar dg@aljabar dg}
	Relatif terhadap struktur monoidal simetris di atas, objek-objek
	$\cate{Alg}\left(\cate{C}(\mathcal{A})\right)$ disebut
	\emph{aljabar diferensial bergradasi} di atas $\mathcal{A}$, disingkat
	\emph{aljabar-dg}; sedangkan objek-objek
	$\cate{CAlg}\left(\cate{C}(\mathcal{A})\right)$ disebut aljabar diferensial
	bergradasi komutatif di atas $\mathcal{A}$, disingkat aljabar-dg komutatif.
\end{definition}

% segment-id: o014.li2.chapter7.differential-graded-structures.p008
Dari sini, untuk suatu aljabar diferensial bergradasi $A$, kita juga dapat
membahas modul-$A$ kiri (atau kanan) $M$, bahkan bimodul; semuanya membentuk
kategori abelian. Dalam keadaan ini, $M$ juga disebut \emph{modul diferensial
bergradasi} di atas $A$, disingkat \emph{modul-dg}. Konsep ini memadukan teori
kompleks dan teori modul serta sangat luas kegunaannya. Latihan-latihan bab
ini akan memberikan pengantar lebih lanjut; pembaca juga dapat merujuk pada
monograf seperti \cite{Yek20}. Untuk menonjolkan struktur diferensial
bergradasinya, selanjutnya kategori modul-dg ditulis $A\dcate{dgMod}$, dan
demikian pula untuk variasi-variasinya.
\index{modul bergradasi diferensial@modul bergradasi diferensial (differential graded module)}
\index{modul dg@modul dg}
\index[sym1]{A-dgMod@$A\dcate{dgMod}$}

% segment-id: o014.li2.chapter7.differential-graded-structures.p009
Funktor pelupa $\cate{C}(\mathcal{A})\to\mathcal{A}^{\Z}$ mempertahankan
$\otimes$, sehingga Proposisi~\ref{prop:lax-monoidal-Alg} menunjukkan bahwa
funktor tersebut mempertahankan struktur aljabar (atau aljabar komutatif).
Dalam kasus $\mathcal{A}=\Bbbk\dcate{Mod}$, mengatakan bahwa $A$ adalah
aljabar-dg sama artinya dengan mengatakan bahwa $A$ sendiri adalah aljabar
bergradasi yang dilengkapi diferensial $d$, dan perkaliannya memenuhi aturan
Leibniz
% segment-id: o014.li2.chapter7.differential-graded-structures.q004
\[ d(xy) = (dx)y + (-1)^a x (dy), \quad x \in A^a, \; y \in A^b. \]
Dengan menyubstitusikan $x=1_A=y$, diperoleh $d(1_A)=0$. Aljabar lawannya
$A^{\opp}$ bersesuaian dengan mengganti perkalian menjadi
$x\mathbin{\cdot^{\opp}}y:=(-1)^{ab}yx$.

% segment-id: o014.li2.chapter7.differential-graded-structures.p010
Serupa dengan itu, bila $\mathcal{A}=\Bbbk\dcate{Mod}$, suatu modul-$A$ kiri
$M$ sama artinya dengan modul-$\Bbbk$ bergradasi $M$ beserta morfisme
perkalian skalar $A\otimes M\to M$ yang memenuhi sifat asosiatif dan
sifat-sifat lainnya, serta dilengkapi diferensial $d$ yang memenuhi aturan
Leibniz
% segment-id: o014.li2.chapter7.differential-graded-structures.q005
\[ d(tm) = (dt)m + (-1)^a t (dm), \quad t \in A^a, \; m \in M^b; \]
untuk modul-$A$ kanan $M$, aturan Leibniz menjadi
% segment-id: o014.li2.chapter7.differential-graded-structures.q006
\[ d(mt) = (dm)t + (-1)^b m (dt), \quad t \in A^a, \; m \in M^b. \]
Untuk $(\mathcal{A},\otimes)$ yang umum, semua sifat di atas dapat dirumuskan
tanpa menyebut unsur. Sebagai contoh, satuan aljabar $A$ bersesuaian dengan
morfisme $\munit\to\Ker\left(d_A^0\right)$ dalam $\mathcal{A}$; lihat
\eqref{eqn:unit-map-into-cplx}.

% segment-id: o014.li2.chapter7.differential-graded-structures.p011
Mengambil kohomologi memberikan funktor aditif
$\Hm=(\Hm^n)_{n\in\Z}:\cate{C}(\mathcal{A})\to\mathcal{A}^{\Z}$. Hasil
berikut menunjukkan bahwa perkalian kompleks menginduksi perkalian pada
kohomologi.

% segment-id: o014.li2.chapter7.differential-graded-structures.pr001
\begin{proposition}\label{prop:Hm-lax-monoidal}
	Funktor $\Hm$ menginduksi
	% segment-id: o014.li2.chapter7.differential-graded-structures.q007
	\[ \cate{Alg}\left(\cate{C}(\mathcal{A})\right) \to \cate{Alg}\left(\mathcal{A}^{\Z}\right),
	\quad \cate{CAlg}\left(\cate{C}(\mathcal{A})\right) \to \cate{CAlg}\left(\mathcal{A}^{\Z}\right). \]
	Jika $A$ adalah aljabar-dg dan $M$ adalah modul-$A$ kiri (atau kanan, atau
	bimodul), maka $\Hm(M)$ adalah modul-$\Hm(A)$ kiri (atau kanan, atau
	bimodul).
\end{proposition}
% segment-id: o014.li2.chapter7.differential-graded-structures.pf001
\begin{proof}
	Menurut Proposisi~\ref{prop:lax-monoidal-Alg}, cukup dibuktikan bahwa $\Hm$
	adalah funktor monoidal longgar. Morfisme kanonik yang diperlukan,
	% segment-id: o014.li2.chapter7.differential-graded-structures.q008
	\[ \xi_{\Hm}(M, N): \Hm(M) \otimes \Hm(N) \to \Hm(M \otimes N), \]
	hampir tautologis: morfisme itu berasal dari morfisme $\kappa$ yang dibahas
	di atas dalam Teorema~\ref{prop:Kunneth-homology} atau
	Proposisi~\ref{prop:bifunctor-cup}. Di sisi lain, kita ambil
	% segment-id: o014.li2.chapter7.differential-graded-structures.q009
	\[ \varphi_{\Hm}: \munit \to \Hm(\munit) = \munit \quad \text{(terkonsentrasi pada derajat nol)} \]
	sebagai $\identity$.
\end{proof}

% segment-id: o014.li2.chapter7.differential-graded-structures.p012
Kita juga dapat mengambil funktor monoidal longgar
$\mathrm{ev}_0:\mathcal{A}^{\Z}\to\mathcal{A}$ dari
Contoh~\ref{eg:ev0-monoidal}. Sekali lagi,
Proposisi~\ref{prop:lax-monoidal-Alg} menunjukkan bahwa funktor ini memetakan
aljabar bergradasi (atau aljabar bergradasi komutatif) ke aljabar (atau
aljabar komutatif), serta memetakan modul ke modul. Perhatikan bahwa
$\mathrm{ev}_0\Hm=\Hm^0$. Dalam penerapan, setiap tahap komposisi
$\cate{C}(\mathcal{A})\xrightarrow{\Hm}\mathcal{A}^{\Z}
\xrightarrow{\mathrm{ev}_0}\mathcal{A}$ menghilangkan informasi; struktur
yang paling kaya tetap berada pada aljabar-dg itu sendiri.

% segment-id: o014.li2.chapter7.differential-graded-structures.eg001
\begin{example}\label{eg:k-linear-A-enrich}
	Misalkan $\Bbbk$ gelanggang komutatif. Untuk sembarang kategori
	$\Bbbk$-linear $\mathcal{A}$, tinjau perkalian pada tingkat kompleks
	$\Hom$ dari \eqref{eqn:Hom-cplx-multiplication}:
	% segment-id: o014.li2.chapter7.differential-graded-structures.q010
	\[ \Hom^a(Y, Z) \times \Hom^b(X, Y) \to \Hom^{a+b}(X, Z), \quad X, Y, Z \in \Obj(\cate{C}(\mathcal{A})). \]

	Karena perkalian memenuhi aturan Leibniz
	(Lema~\ref{prop:Hom-cplx-Leibniz}), perkalian itu menentukan morfisme dalam
	$\cate{C}(\Bbbk\dcate{Mod})$
	% segment-id: o014.li2.chapter7.differential-graded-structures.q011
	\[ \Hom^\bullet(Y, Z) \otimes \Hom^\bullet(X, Y) \to \Hom^\bullet(X, Z). \]
	Selain itu, perkalian tersebut asosiatif, sehingga
	$\End^\bullet(X):=\Hom^\bullet(X,X)$ menjadi aljabar di atas
	$(\cate{C}(\Bbbk\dcate{Mod}),\otimes)$ dengan satuan
	$\identity_X\in\End^0(X)$. Lebih lanjut, $\Hom^\bullet(X,Y)$ menjadi
	bimodul-$(\End^\bullet(Y),\End^\bullet(X))$.

	Mengambil kohomologi $\Hm=(\Hm^n)_n$ dari $\Hom^\bullet(X,Y)$ hanya
	menghasilkan
	% segment-id: o014.li2.chapter7.differential-graded-structures.q012
	\[ \bigoplus_{n \in \Z} \Hom_{\cate{K}(\mathcal{A})}(X, Y[n]); \]
	sedangkan dengan menerapkan \eqref{eqn:unit-map-into-cplx} pada
	$\Hom^\bullet(X,Y)$, tampak bahwa $\Hom_{\cate{C}(\mathcal{A})}$ dapat
	direkonstruksi melalui
	% segment-id: o014.li2.chapter7.differential-graded-structures.q013
	\[ \Hom_{\cate{C}(\Bbbk\dcate{Mod})}\left( \Bbbk, \Hom^\bullet(X, Y) \right) \simeq \Hom_{\cate{C}(\mathcal{A})}(X, Y)  \]
	dengan $\Bbbk$ berperan sebagai objek satuan dari
	$(\cate{C}(\Bbbk\dcate{Mod}),\otimes)$.
\end{example}

% segment-id: o014.li2.chapter7.differential-graded-structures.p013
Jadi, himpunan $\Hom_{\cate{C}(\mathcal{A})}$ dan komposisi morfisme dalam
$\cate{C}(\mathcal{A})$ hanyalah bayangan kompleks $\Hom$ beserta
perkaliannya, yang diperoleh dengan mengambil
$\Hom_{\cate{C}(\Bbbk\dcate{Mod})}(\Bbbk,\mathord\cdot)$.

\index{kategori!diperkaya@diperkaya (enriched)}
\index[sym1]{Hom-internal@$\iHom$}
% segment-id: o014.li2.chapter7.differential-graded-structures.p014
Berangkat dari contoh $\cate{C}(\mathcal{A})$, kita secara alami sampai pada
konsep kategori diferensial bergradasi. Konsep ini melibatkan
\emph{kategori yang diperkaya-$\mathcal{V}$} yang diperkenalkan dalam
\cite[Definisi 3.4.1]{Li1}. Singkatnya, untuk sembarang kategori monoidal
$\mathcal{V}$, sebuah kategori yang diperkaya-$\mathcal{V}$, $\mathcal{C}$,
tetap terdiri atas objek dan morfisme, tetapi himpunan $\Hom$ diganti dengan
objek $\Hom$ dalam $\mathcal{V}$, yang ditulis
$\iHom=\iHom_{\mathcal{C}}$ untuk membedakannya. Komposisi morfisme diganti
dengan morfisme dalam $\mathcal{V}$
% segment-id: o014.li2.chapter7.differential-graded-structures.q014
\[ \iHom(Y, Z) \otimes \iHom(X, Y) \to \iHom(X, Z), \]
sedangkan morfisme identitas $\identity_X\in\Hom(X,X)$ diganti dengan
$\identity_X:\munit\to\iHom(X,X)$. Semua ini adalah morfisme dalam
$\mathcal{V}$, dan sifat-sifat yang berkaitan dengannya dinyatakan melalui
diagram komutatif dalam $\mathcal{V}$.

% segment-id: o014.li2.chapter7.differential-graded-structures.p015
Untuk versi diperkaya dari funktor dan morfisme di antaranya, lihat
\cite[Definisi 3.4.1, 3.4.2, 3.4.4]{Li1}. Semuanya didefinisikan dengan cara
yang sewajarnya; singkatnya, versi diperkaya dari funktor mengangkat
pemetaan antarhimpunan $\Hom$ menjadi morfisme antarobjek $\iHom$.

% segment-id: o014.li2.chapter7.differential-graded-structures.p016
Sebagai contoh, kategori-$\Bbbk\dcate{Mod}$ tidak lain adalah kategori yang
diperkaya atas $(\Bbbk\dcate{Mod},\otimes)$.

% segment-id: o014.li2.chapter7.differential-graded-structures.p017
Sebaliknya, untuk sementara kita sebut kategori yang tidak diperkaya sebagai
kategori biasa. Jika persoalan ukuran himpunan diabaikan, kategori biasa sama
artinya dengan kategori yang diperkaya atas $(\cate{Set},\times)$.

% segment-id: o014.li2.chapter7.differential-graded-structures.d002
\begin{definition}\label{def:dgCat}
	\index{kategori!bergradasi diferensial@bergradasi diferensial (differential graded)}
	\index{kategori dg@kategori dg}
	Misalkan $\Bbbk$ gelanggang komutatif. Tinjau kategori monoidal simetris
	$(\cate{C}(\Bbbk\dcate{Mod}),\otimes)$ yang bersesuaian. Kategori yang
	diperkaya atas $(\cate{C}(\Bbbk\dcate{Mod}),\otimes)$ disebut
	\emph{kategori diferensial bergradasi} di atas $\Bbbk$, disingkat
	\emph{kategori-dg} di atas $\Bbbk$. Funktor antarkategori tersebut
	dan morfisme antarfunktor didefinisikan seperti dalam kategori yang
	diperkaya\footnote{Banyak pustaka mengenai kategori tak hingga menyebut
	DG-kategori serta kategori $\cate{DGCat}$ yang dibentuk olehnya, dan
	kadang notasinya juga bertumpang tindih dengan notasi buku ini. Walaupun
	DG-kategori tersebut berkaitan erat dengan kategori-dg yang dibahas di
	sini, keduanya bukan hal yang sama.}.
\end{definition}

% segment-id: o014.li2.chapter7.differential-graded-structures.p018
Menurut definisi, sebuah objek $X$ dalam kategori-dg memberikan aljabar-dg
$\End^\bullet(X)$.

% segment-id: o014.li2.chapter7.differential-graded-structures.eg002
\begin{example}\label{eg:CA-dg-cat}
	Contoh~\ref{eg:k-linear-A-enrich} di atas mengatakan bahwa
	$\cate{C}(\mathcal{A})$ secara alami menjadi kategori-dg di atas $\Bbbk$,
	dengan $\mathcal{A}$ sembarang kategori $\Bbbk$-linear. Sebagai contoh,
	untuk setiap aljabar-$\Bbbk$ $R$, kategori kompleks modul-$R$ kiri (atau
	kanan) secara otomatis merupakan kategori-dg di atas $\Bbbk$.
\end{example}

% segment-id: o014.li2.chapter7.differential-graded-structures.eg003
\begin{example}\label{eg:dg-Hom-cplx}
	\index{Hom fuxing}
	\index[sym1]{Hom-A-bullet@$\Hom^{\bullet}_A$}
	Misalkan $\Bbbk$ gelanggang komutatif dan $A$ aljabar-dg di atas
	$\mathcal{A}:=\Bbbk\dcate{Mod}$. Untuk modul-$A$ kiri $M$ dan $N$,
	definisikan subkompleks dari $\Hom^\bullet_{\Bbbk}(M,N)$ sebagai berikut:
	% segment-id: o014.li2.chapter7.differential-graded-structures.q015
	\[ \Hom^p_A(M, N) := \left\{\begin{array}{r|l}
		f \in \Hom^p_{\Bbbk}(M, N) & \forall m \in M, \; \forall k \in \Z, \; \forall t \in A^k, \\
		& f(tm) = (-1)^{pk} tf(m)
	\end{array}\right\}. \]
	Untuk modul-$A$ kiri $L$, $M$, dan $N$ yang diberikan, operasi komposisi
	pada kompleks $\Hom^\bullet_{\Bbbk}$ menginduksi morfisme antarkompleks
	% segment-id: o014.li2.chapter7.differential-graded-structures.q016
	\[ \Hom^p_A(M, N) \otimes \Hom^q_A(L, M) \to \Hom^{p+q}_A(L, N). \]
	Dengan demikian, kategori modul-$A$ kiri $A\dcate{dgMod}$ diperkaya menjadi
	kategori-dg. Untuk modul-$A$ kanan, syarat dalam definisi
	$\Hom^p_A(M,N)$ harus diganti dengan $f(mt)=f(m)t$ (renungkanlah).
	Kasus bimodul diperlakukan dengan cara serupa.

	Definisi di atas juga dapat dirumuskan tanpa bergantung pada unsur
	himpunan, sehingga meluas ke $\mathcal{A}$ yang umum.
\end{example}

% segment-id: o014.li2.chapter7.differential-graded-structures.p019
Lebih tepatnya, kategori-dg memperkaya semua himpunan $\Hom$ dalam kategori
biasa menjadi kompleks $\Hom^\bullet$, sedangkan funktor dan morfisme versi
dg harus dinyatakan di dalam $\cate{C}(\Bbbk\dcate{Mod})$. Secara terperinci,
% segment-id: o014.li2.chapter7.differential-graded-structures.l001
\begin{itemize}
	\item Sebuah \emph{funktor-dg} $F:\mathcal{C}\to\mathcal{D}$ antara
	kategori-dg ditentukan oleh pemetaan antarhimpunan objek
	$F:\Obj(\mathcal{C})\to\Obj(\mathcal{D})$ dan morfisme antarkompleks
	$\Hom$
	$\Hom^\bullet_{\mathcal{C}}(X,Y)\to
	\Hom^\bullet_{\mathcal{D}}(FX,FY)$. Morfisme yang terakhir adalah morfisme
	dalam $\cate{C}(\Bbbk\dcate{Mod})$ dan disyaratkan membuat kedua diagram
	berikut komutatif:
	% segment-id: o014.li2.chapter7.differential-graded-structures.q017
	\begin{equation*}\begin{gathered}
		\begin{tikzcd}[column sep=small]
			\Hom_{\mathcal{C}}^\bullet(Y, Z) \otimes \Hom_{\mathcal{C}}^\bullet(X, Y) \arrow[r] \arrow[d] & \Hom_{\mathcal{C}}^\bullet(X, Z) \arrow[d] \\
			\Hom_{\mathcal{D}}^\bullet(FY, FZ) \otimes \Hom_{\mathcal{D}}^\bullet(FX, FY) \arrow[r] & \Hom_{\mathcal{D}}^\bullet(FX, FZ)
		\end{tikzcd} \\
		\begin{tikzcd}[column sep=small]
			\Bbbk \arrow[r] \arrow[rd] & \Hom_{\mathcal{C}}^\bullet(X, X) \arrow[d] \\
			& \Hom_{\mathcal{D}}^\bullet(FX, FX) ;
		\end{tikzcd}
	\end{gathered}\end{equation*}
	\index{funktor dg@funktor dg}
	\item Dalam pengertian klasik, suatu morfisme, atau transformasi alami,
	$\phi$ antara funktor-dg $F,G:\mathcal{C}\to\mathcal{D}$ adalah keluarga
	unsur
	% segment-id: o014.li2.chapter7.differential-graded-structures.q018
	\begin{multline*}
		\phi_X \in \Ker\left[ \Hom_{\mathcal{D}}^0(FX, GX) \to \Hom_{\mathcal{D}}^1(FX, GX) \right] \\
		\xlongequal{\text{\eqref{eqn:unit-map-into-cplx}}} \Hom_{\cate{C}(\Bbbk\dcate{Mod})}\left(\Bbbk, \Hom_{\mathcal{D}}^\bullet(FX, GX) \right) ,
	\end{multline*}
	dengan $X$ merentang $\Obj(\mathcal{C})$, sedemikian sehingga untuk
	semua $m\in\Z$ dan $f\in\Hom_{\mathcal{C}}^m(X,Y)$ berlaku
	% segment-id: o014.li2.chapter7.differential-graded-structures.q019
	\[ (Gf) \phi_X = \phi_Y (Ff) \; \in \Hom_{\mathcal{D}}^m(FX, GY), \]
	dengan perkalian dipahami sebagai perkalian pada kompleks $\Hom$.
\end{itemize}
Dengan demikian, kita dapat membicarakan isomorfisme, ekuivalensi, adjoin,
dan konsep-konsep lain untuk kategori-dg. Unsur dari
$\Hom^n_{\mathcal{C}}(X,Y)$ kadang-kadang disebut \emph{morfisme berderajat
$n$} dari $X$ ke $Y$.

% segment-id: o014.li2.chapter7.differential-graded-structures.p020
Untuk kategori monoidal umum $\mathcal{V}$, kita dapat memandang
$\Hom_{\mathcal{V}}(\munit,M)$ sebagai himpunan ``unsur'' dari
$M\in\Obj(\mathcal{V})$; setidaknya pandangan ini wajar ketika
$\mathcal{V}=\cate{Set}$ atau $\Bbbk\dcate{Mod}$. Mengikuti gagasan ini,
setiap kategori yang diperkaya-$\mathcal{V}$, $\mathcal{C}$, dapat diturunkan
menjadi kategori biasa dengan menetapkan
% segment-id: o014.li2.chapter7.differential-graded-structures.q020
\[ \Hom_{\mathcal{C}}(X, Y) := \Hom_{\mathcal{V}}\left( \munit, \iHom_{\mathcal{C}}(X, Y) \right), \]
dan kategori biasa ini tetap ditulis $\mathcal{C}$; lihat
\cite[Catatan 3.4.3]{Li1} untuk perinciannya. Dalam contoh kategori-dg
$\cate{C}(\mathcal{A})$, hasil penurunan tersebut adalah kategori kompleks
biasa $\cate{C}(\mathcal{A})$. Jadi, notasi kita konsisten.

% segment-id: o014.li2.chapter7.differential-graded-structures.p021
Operasi lain berasal dari funktor monoidal longgar. Berikut pernyataan versi
umumnya.

% segment-id: o014.li2.chapter7.differential-graded-structures.pr002
\begin{proposition}\label{prop:lax-monoidal-enriched}
	Misalkan $F:\mathcal{V}\to\mathcal{V}'$ funktor monoidal longgar antara
	kategori-kategori monoidal. Untuk setiap kategori yang
	diperkaya-$\mathcal{V}$, $\mathcal{C}$, definisikan kategori yang
	diperkaya-$\mathcal{V}'$, $F(\mathcal{C})$, sebagai berikut: tetapkan
	$\Obj(F(\mathcal{C}))=\Obj(\mathcal{C})$, dan untuk setiap objek $X,Y$,
	tetapkan
	% segment-id: o014.li2.chapter7.differential-graded-structures.q021
	\begin{gather*}
		\iHom_{F(\mathcal{C})}(X, Y) := F\iHom_{\mathcal{C}}(X, Y), \\
		\identity_X: \munit_{\mathcal{V}'} \xrightarrow{\varphi_F} F(\munit_{\mathcal{V}}) \to F\iHom_{\mathcal{C}}(X, X).
	\end{gather*}
	Komposisi morfisme ditentukan oleh
	% segment-id: o014.li2.chapter7.differential-graded-structures.q022
	\begin{multline*}
		F\iHom_{\mathcal{C}}(Y, Z) \otimes F\iHom_{\mathcal{C}}(X, Y) \\
		\xrightarrow{\xi_F} F\left( \iHom_{\mathcal{C}}(Y, Z) \otimes \iHom_{\mathcal{C}}(X, Y) \right)
		\to F\iHom_{\mathcal{C}}(X, Z),
	\end{multline*}
	dan pada kategori biasa yang bersesuaian, konstruksi ini memberikan
	funktor dalam pengertian biasa $\mathcal{C}\to F(\mathcal{C})$.
\end{proposition}
% segment-id: o014.li2.chapter7.differential-graded-structures.pf002
\begin{proof}
	Struktur funktor monoidal longgar $(F,\xi_F,\varphi_F)$ menyediakan semua
	sifat yang diperlukan oleh komposisi morfisme. Pada tingkat kategori biasa,
	funktor tersebut adalah identitas pada objek dan, pada himpunan morfisme,
	memetakan setiap $f:\munit_{\mathcal{V}}\to
	\iHom_{\mathcal{C}}(X,Y)$ ke komposisi
	$\munit_{\mathcal{V}'}\xrightarrow{\varphi_F}F(\munit_{\mathcal{V}})
	\xrightarrow{Ff}F\iHom_{\mathcal{C}}(X,Y)$.
\end{proof}

% segment-id: o014.li2.chapter7.differential-graded-structures.p022
Dengan membandingkan konstruksi ini dengan konstruksi dalam
Proposisi~\ref{prop:lax-monoidal-Alg}, terlihat bahwa
$\iHom_{F(\mathcal{C})}(X,X)$ sebagai objek $\cate{Alg}(\mathcal{V})$
diinduksi dari $\iHom_{\mathcal{C}}(X,X)$; demikian pula struktur bimodul
pada $\iHom$.

% segment-id: o014.li2.chapter7.differential-graded-structures.r001
\begin{remark}\label{rem:dg-homotopy-cat}
	\index{kategori!homotopi@homotopi (homotopy)}
	Proposisi~\ref{prop:lax-monoidal-enriched} memberikan cara lain untuk
	beralih dari kategori-dg ke kategori biasa. Terapkan funktor monoidal
	longgar $\Hm$ dari Proposisi~\ref{prop:Hm-lax-monoidal} pada sembarang
	kategori-dg $\mathcal{C}$ di atas $\Bbbk$. Hasilnya adalah kategori yang
	diperkaya atas $(\Bbbk\dcate{Mod}^{\Z},\otimes)$, yaitu
	$\Hm(\mathcal{C})$, yang memenuhi
	% segment-id: o014.li2.chapter7.differential-graded-structures.q023
	\[ \Obj(\Hm(\mathcal{C})) = \Obj(\mathcal{C}), \quad \iHom_{\Hm(\mathcal{C})}(X, Y) = \left( \Hm^n \Hom^\bullet(X, Y) \right)_{n \in \Z}. \]
	Dengan mengambil bagian $n=0$, atau, dengan kata lain, dengan menerapkan
	funktor monoidal longgar $\Hm^0$ pada $\mathcal{C}$, kita memperoleh
	kategori-$\Bbbk\dcate{Mod}$ biasa yang disebut \emph{kategori homotopi}
	$\mathrm{h}(\mathcal{C})$ dari $\mathcal{C}$. Kategori itu memenuhi
	% segment-id: o014.li2.chapter7.differential-graded-structures.q024
	\[ \Obj(\mathrm{h}\mathcal{C}) = \Obj(\mathcal{C}), \quad \Hom_{\mathrm{h}\mathcal{C}}(X, Y) = \Hm^0 \Hom^\bullet(X, Y). \]
	Dalam kasus khusus $\mathcal{C}=\cate{C}(\mathcal{A})$, dengan
	menyubstitusikan Definisi~\ref{def:cplx-homotopy-cat}, langsung diperoleh
	% segment-id: o014.li2.chapter7.differential-graded-structures.q025
	\[ \mathrm{h}(\cate{C}(\mathcal{A})) = \cate{K}(\mathcal{A}). \]
\end{remark}

% segment-id: o014.li2.chapter7.differential-graded-structures.p023
Seperti terlihat di atas, dunia kategori-dg mengizinkan berbagai operasi
yang bernuansa homotopi. Maknanya pada akhirnya harus dijelaskan melalui
penerapan, dan juga akan melibatkan bahasa abstrak kategori model. Berikut
ini hanya diberikan garis besar konsep-konsep sederhana tanpa pengembangan
lebih lanjut.

% segment-id: o014.li2.chapter7.differential-graded-structures.d003
\begin{definition}
	Misalkan $F:\mathcal{C}\to\mathcal{D}$ funktor-dg antara kategori-dg. Jika
	% segment-id: o014.li2.chapter7.differential-graded-structures.q026
	\[ \Hom^\bullet_{\mathcal{C}}(X, Y) \to \Hom^\bullet_{\mathcal{D}}(FX, FY) \]
	merupakan kuasi-isomorfisme untuk semua
	$X,Y\in\Obj(\mathcal{C})$, maka $F$ disebut \emph{kuasi-penuh-setia}; jika
	funktor terinduksi
	$\mathrm{h}F:\mathrm{h}\mathcal{C}\to\mathrm{h}\mathcal{D}$ esensial
	surjektif, maka $F$ disebut \emph{kuasi-esensial-surjektif}. Funktor-dg
	yang sekaligus kuasi-penuh-setia dan kuasi-esensial-surjektif disebut
	\emph{kuasi-ekuivalensi}.
\end{definition}

% segment-id: o014.li2.chapter7.differential-graded-structures.p024
Karena itu, kategori-dg $\mathcal{C}$ dan $\mathcal{D}$ yang kuasi-ekuivalen
mempunyai kategori homotopi yang ekuivalen.

% segment-id: o014.li2.chapter7.differential-graded-structures.p025
Kita akan melanjutkan pembahasan ringkas mengenai kategori-dg dalam
\S\ref{sec:dg-closed}.
